%!TeX program = xelatex
\documentclass{SYSUReport}
\usepackage{amsmath}
\usepackage{tikz}
\usetikzlibrary{arrows.meta, positioning}
\usepackage{adjustbox}
\usepackage{booktabs}
\usepackage{multirow}
\usepackage{graphicx}
\usepackage{longtable}
\usepackage{listings}
\usepackage{xcolor}
\usepackage{url}
\urlstyle{tt}
\def\UrlBreaks{\do\/\do\.\do\-\do\_}

\lstdefinestyle{mystyle}{
    backgroundcolor=\color{gray!10},
    basicstyle=\ttfamily\small,
    keywordstyle=\color{blue!70!black},
    commentstyle=\color{green!40!black},
    stringstyle=\color{purple!60!black},
    numbers=left,
    numberstyle=\tiny\color{gray},
    frame=single,
    breaklines=true,
    captionpos=b,
    showstringspaces=false
}
\lstdefinestyle{shell}{
    style=mystyle,
    numbers=none,
    language=bash
}
\lstset{style=mystyle}

\headl{学号 姓名}
\headc{}
\headr{数据安全和隐私保护}
\lessonTitle{数据安全和隐私保护}
\reportTitle{数据安全和隐私保护第一次实验}
\stuname{姓名}
\stuid{学号}
\inst{网络空间安全学院}
\major{网络空间安全}
\date{2026 年 9 月 10 日}

\begin{document}

\cover
\thispagestyle{empty}
\clearpage

\tableofcontents
\clearpage

\pagenumbering{arabic}
\setcounter{page}{1}

\section{实验要求}

本次实验围绕科研过程中的基础工具展开，包含三个部分。

\begin{enumerate}
    \item LaTeX 图片排版：掌握 \verb|\includegraphics| 的使用方法，能够设置图片尺寸，为图片配置图编号、图标题并在正文中引用；优先使用 PDF、SVG、EPS 等矢量格式，避免使用低分辨率截图作为结果图。报告中至少插入一张实验结果图。
    \item Git 与 GitHub：完成 Git 的安装与配置，完成仓库创建、克隆、提交、上传与更新。
    \item Codex 与 CC Switch：完成 Codex 的安装与基本配置，安装 CC Switch 并完成供应商配置与切换，使用 Codex 完成至少一次任务。
\end{enumerate}

按实验说明，Git、Codex、CC Switch 已在此前安装并配置完成的情况下可跳过安装步骤。本报告对应的三项工具均已在实验前完成安装与配置，因此第 \ref{sec:process} 节记录的是配置状态核对与实际操作过程。报告中涉及 API Key、Endpoint 凭据的位置均已脱敏。

\section{实验原理}

\subsection{位图与矢量图}

位图由像素矩阵组成，放大后受采样密度限制，边缘会出现锯齿或模糊。矢量图记录曲线、文字和几何对象的描述，渲染时按输出分辨率重新光栅化，缩放后仍保持清晰。折线图、柱状图、网络结构图等由绘图程序生成的图形，可直接导出为 PDF 或 SVG，避免经过截图带来的清晰度损失。

LaTeX 通过 \verb|graphicx| 宏包将图片按版面尺寸排入文档。\verb|figure| 环境提供浮动定位，\verb|\caption| 生成图题与图编号，\verb|\label| 与 \verb|\ref| 建立正文引用关系。

\subsection{Git 的三个区域}

Git 是分布式版本控制系统。工作区保存正在编辑的文件，暂存区保存准备提交的变化，\verb|commit| 将暂存内容记录到本地仓库，\verb|push| 将本地提交上传到远程仓库，\verb|pull| 取得远程提交并合并到当前分支。\verb|status|、\verb|diff|、\verb|log| 用于在提交前后核对状态、内容差异和历史记录。数据流向如图 \ref{fig:gitflow} 所示。

\begin{figure}[htbp]
    \centering
    \begin{tikzpicture}[
        node distance=8mm,
        box/.style={draw=myteal, fill=myteal!6, rounded corners=2pt,
                    minimum width=25mm, minimum height=11mm, align=center, font=\small},
        arr/.style={-{Stealth[length=2.5mm]}, thick, draw=myteal!80!black, font=\scriptsize}
    ]
        \node[box] (work) {工作区};
        \node[box, right=16mm of work] (stage) {暂存区};
        \node[box, right=16mm of stage] (local) {本地仓库};
        \node[box, right=16mm of local] (remote) {GitHub\\远程仓库};

        \draw[arr] (work) -- node[above]{git add} (stage);
        \draw[arr] (stage) -- node[above]{git commit} (local);
        \draw[arr] (local) -- node[above]{git push} (remote);
        \draw[arr] (remote.south) .. controls +(0,-9mm) and +(0,-9mm) ..
              node[below]{git pull} (work.south);
    \end{tikzpicture}
    \caption{Git 工作区、暂存区、本地仓库与远程仓库之间的数据流向}
    \label{fig:gitflow}
\end{figure}

\subsection{Codex 与 CC Switch}

Codex 在用户授权范围内读取项目、修改文件并运行本地命令，根据自然语言需求和仓库上下文制定并执行操作。生成结果可能存在错误、遗漏或不安全修改，因此流程中保留权限确认、命令审阅、程序运行、测试和 Git 差异检查等核对环节。

CC Switch 将不同供应商的端点、模型和认证配置保存为可切换条目。切换时更新 Codex 的配置文件 \path{~/.codex/config.toml}，命令行工具重启后生效。使用第三方供应商时，密钥保护、端口暴露和服务条款均需关注。

\section{实验设置}

实验在本地 macOS 主机上完成工具操作，LaTeX 报告在远程 Linux 服务器上使用 XeLaTeX 编译。软件版本见表 \ref{tab:env}。

\begin{table}[htbp]
    \centering
    \caption{实验环境}
    \label{tab:env}
    \begin{tabular}{ll}
        \toprule
        项目 & 版本或配置 \\
        \midrule
        操作系统（工具操作）  & macOS，Darwin 27.0.0，arm64 \\
        操作系统（报告编译）  & Linux 服务器，通过 SSH 访问 \\
        Git                 & 2.50.1 (Apple Git-155) \\
        Python              & 3.9，matplotlib 3.9.4 \\
        Codex CLI           & 0.153.4 \\
        CC Switch           & 3.14.1（Homebrew Cask 安装） \\
        LaTeX               & TeX Live，XeLaTeX 引擎 \\
        报告模板            & SYSUReport.cls \\
        \bottomrule
    \end{tabular}
\end{table}

实验仓库地址为 \url{https://github.com/Candlest/research-tools-week1}，目录结构为 \verb|code/|、\verb|result/|、\verb|report/|、\verb|figures/| 四类。

\section{实验过程}
\label{sec:process}

\subsection{LaTeX 图片排版}

实验结果图由 \verb|code/plot_word_freq.py| 生成。该程序读取 \verb|result/word_freq.csv|，使用 matplotlib 的 Agg 后端绘制横向柱状图，并直接导出为 PDF 矢量文件，未经过截图环节。关键设置为 \verb|pdf.fonttype=42|，使文字以可缩放字形嵌入。

\begin{lstlisting}[style=shell, caption={生成矢量结果图}]
$ python3 code/plot_word_freq.py
wrote figures/word_freq.pdf

$ file figures/word_freq.pdf
figures/word_freq.pdf: PDF document, version 1.4, 1 pages
\end{lstlisting}

在导言区加载 \verb|graphicx| 宏包后，按下列方式插入图片。\verb|width| 以 \verb|\textwidth| 为基准设置尺寸，\verb|\caption| 与 \verb|\label| 分别提供图题和引用标签。

\begin{lstlisting}[language=TeX, caption={图片插入代码}]
\usepackage{graphicx}
\begin{figure}[htbp]
  \centering
  \includegraphics[width=0.80\textwidth]{figures/word_freq.pdf}
  \caption{sample.txt 词频统计结果（前 10 个词）}
  \label{fig:wordfreq}
\end{figure}
\end{lstlisting}

插入结果如图 \ref{fig:wordfreq} 所示。图中横轴为出现次数，纵轴为词。将 PDF 页面放大后，坐标轴标签、刻度数字和数值标注保持清晰。

\begin{figure}[htbp]
    \centering
    \includegraphics[width=0.80\textwidth]{figures/word_freq.pdf}
    \caption{sample.txt 词频统计结果（前 10 个词）}
    \label{fig:wordfreq}
\end{figure}

\subsection{Git 与 GitHub}

\subsubsection{配置核对}

Git 已在实验前安装，提交身份已配置。核对结果如下，配置文件中不含密码、令牌或代理设置。

\begin{lstlisting}[style=shell, caption={Git 版本与用户级配置}]
$ git --version
git version 2.50.1 (Apple Git-155)

$ git config --global --list
user.email=candlest1112@gmail.com
user.name=Candlest
http.version=HTTP/1.1
alias.dotfiles=!git --git-dir=$HOME/.dotfiles --work-tree=$HOME
\end{lstlisting}

\subsubsection{创建仓库并克隆}

在 GitHub 创建公开仓库 \verb|research-tools-week1| 并初始化 README，随后克隆到本地并建立 \verb|code/|、\verb|result/|、\verb|report/|、\verb|figures/| 四个目录。空目录不会被 Git 记录，因此在 \verb|result/| 和 \verb|report/| 中各放置一个说明文件。

\begin{lstlisting}[style=shell, caption={创建并克隆仓库}]
$ gh repo create research-tools-week1 --public --add-readme
https://github.com/Candlest/research-tools-week1

$ git clone https://github.com/Candlest/research-tools-week1.git
Cloning into 'research-tools-week1'...

$ cd research-tools-week1
\end{lstlisting}

\subsubsection{第一次提交}

提交前依次执行 \verb|status| 和 \verb|diff|，确认变化范围与内容，再明确列出待暂存文件，未使用 \verb|git add .|。

\begin{lstlisting}[style=shell, caption={检查修改并创建第一次提交}]
$ git status
On branch main
Your branch is up to date with 'origin/main'.

Changes not staged for commit:
	modified:   README.md

Untracked files:
	code/
	report/
	result/

$ git diff
diff --git a/README.md b/README.md
index 16bd6b3..4e73d8d 100644
--- a/README.md
+++ b/README.md
@@ -1,2 +1,19 @@
 # research-tools-week1
-...
+...

$ git add README.md code/demo.py result/README.md report/README.md

$ git diff --staged --stat
 README.md        | 19 ++++++++++++++++++-
 code/demo.py     | 21 +++++++++++++++++++++
 report/README.md |  1 +
 result/README.md |  1 +
 4 files changed, 41 insertions(+), 1 deletion(-)

$ git commit -m "add project layout and first runnable demo program"
[main 092cb4d] add project layout and first runnable demo program
 4 files changed, 41 insertions(+), 1 deletion(-)
\end{lstlisting}

\subsubsection{上传到远程仓库}

首次向 main 分支上传时使用 \verb|-u| 建立上游关系。

\begin{lstlisting}[style=shell, caption={建立上游关系并上传}]
$ git branch --show-current
main

$ git push -u origin main
To https://github.com/Candlest/research-tools-week1.git
   03db75d..092cb4d  main -> main
branch 'main' set up to track 'origin/main'.

$ git log --oneline --graph --decorate -n 10
* 092cb4d (HEAD -> main, origin/main, origin/HEAD) add project layout and first runnable demo program
* 03db75d Initial commit
\end{lstlisting}

\subsubsection{第二次提交}

第二次提交加入 Codex 生成的词频统计程序、单元测试、绘图程序、样例输入、统计结果 CSV 和结果图 PDF。

\begin{lstlisting}[style=shell, caption={第二次提交并上传}]
$ git add code/text_stats.py code/test_text_stats.py code/plot_word_freq.py \
        result/sample.txt result/word_freq.csv figures/word_freq.pdf

$ git diff --staged --stat
 code/plot_word_freq.py  |  44 ++++++++++++++++++
 code/test_text_stats.py |  61 +++++++++++++++++++++++
 code/text_stats.py      |  35 ++++++++++++
 figures/word_freq.pdf   | Bin 0 -> 10824 bytes
 result/sample.txt       |  16 ++++++
 result/word_freq.csv    |   9 +++++
 6 files changed, 165 insertions(+)

$ git commit -m "add word frequency program, unit tests and result figure"
[main d6f5db7] add word frequency program, unit tests and result figure
 6 files changed, 165 insertions(+)

$ git push
To https://github.com/Candlest/research-tools-week1.git
   092cb4d..d6f5db7  main -> main
\end{lstlisting}

\subsubsection{同步远程更新}

在 GitHub 端修改 README 中的一行文字并提交，随后在本地确认工作区干净并执行 \verb|git pull|。本次为快进合并，未出现冲突。

\begin{lstlisting}[style=shell, caption={拉取远程更新}]
$ git status
On branch main
Your branch is up to date with 'origin/main'.

nothing to commit, working tree clean

$ git pull
From https://github.com/Candlest/research-tools-week1
   d6f5db7..457f9bd  main       -> origin/main
Updating d6f5db7..457f9bd
Fast-forward
 README.md | 2 +-
 1 file changed, 1 insertion(+), 1 deletion(-)

$ git log --oneline --graph --decorate -n 10
* 457f9bd (HEAD -> main, origin/main, origin/HEAD) edit README line on GitHub web
* d6f5db7 add word frequency program, unit tests and result figure
* 092cb4d add project layout and first runnable demo program
* 03db75d Initial commit
\end{lstlisting}

历史记录中包含两次本地提交和一次远程更新，与实验要求一致。

\subsection{Codex 辅助编程}

\subsubsection{版本与配置核对}

Codex CLI 已在实验前安装。当前通过 API Key 方式使用第三方供应商，因此 \verb|codex login status| 显示未使用 ChatGPT 账号登录。

\begin{lstlisting}[style=shell, caption={Codex 版本与登录状态}]
$ codex --version
codex-cli 0.153.4

$ codex login status
Not logged in
\end{lstlisting}

用户级配置位于 \path{~/.codex/config.toml}，关键字段如下，其中 API Key 由 CC Switch 写入 \verb|auth.json| 并已脱敏。

\begin{lstlisting}[caption={config.toml 关键字段（脱敏）}]
model_provider = "custom"
model = "gpt-6-astra"
model_reasoning_effort = "low"

[model_providers.custom]
name = "<REDACTED>"
base_url = "<REDACTED>"
wire_api = "responses"
\end{lstlisting}

\subsubsection{只读阅读项目}

第一步限制 Codex 为只读沙盒，要求其阅读项目并说明结构，不修改文件。

\begin{lstlisting}[style=shell, caption={要求 Codex 只读阅读项目}]
$ codex exec --sandbox read-only \
    "read the project, describe the layout, entry point and how to run it, do not modify files"
\end{lstlisting}

Codex 输出的目录结构与入口说明如下，与仓库实际内容一致。它同时指出 README 中提到的 \verb|code/text_stats.py| 与 \verb|result/sample.txt| 当时尚不存在，这一判断正确。

\begin{lstlisting}[caption={Codex 阅读结果}]
research-tools-week1/
|-- README.md          # project description and run commands
|-- code/
|   `-- demo.py        # character and line counting
|-- result/README.md   # results and figures
|-- report/README.md   # LaTeX report
`-- figures/           # vector figures for the report

Entry point: main() in code/demo.py. Standard library only, Python 3.9+.
\end{lstlisting}

\subsubsection{生成程序与测试}

第二步在 workspace-write 沙盒下要求 Codex 编写词频统计程序、样例输入和单元测试。任务描述明确了切分规则、排序方式、参数形式和输出文件。

\begin{lstlisting}[style=shell, caption={要求 Codex 完成编程任务}]
$ codex exec --sandbox workspace-write \
    "add code/text_stats.py: read a text file, count English word frequency
     (lowercase, split on non-letters, drop empty tokens), print the top N words
     in descending order with --top (default 10), and write result/word_freq.csv
     with header word,count. Also add result/sample.txt (English, 15+ lines) and
     code/test_text_stats.py with at least 3 unittest cases. Run the tests.
     Standard library only."
\end{lstlisting}

Codex 新增了三个文件并自行运行了测试。退出 Codex 后使用 Git 独立核对修改范围，未发现预期外的文件变动。

\begin{lstlisting}[style=shell, caption={核对修改并运行程序与测试}]
$ git status --short
?? code/test_text_stats.py
?? code/text_stats.py
?? result/sample.txt
?? result/word_freq.csv

$ python3 code/text_stats.py result/sample.txt --top 10
research	5
text	5
a	4
words	4
analysis	3
tools	3
file	3
each	3
in	3
count	3

$ python3 -B -m unittest discover -s code -p 'test_*.py' -v
test_cli_default_and_custom_top_with_csv ... ok
test_descending_frequency_and_top_limit ... ok
test_empty_tokens_are_filtered ... ok
test_lowercase_and_nonletter_separators ... ok

----------------------------------------------------------------------
Ran 4 tests in 0.089s

OK
\end{lstlisting}

4 个测试全部通过，覆盖大小写转换与非字母切分、空串过滤、降序排序与数量限制、命令行参数与 CSV 内容。生成的 \verb|result/word_freq.csv| 作为图 \ref{fig:wordfreq} 的数据来源。核对无误后执行第二次提交并上传。

\subsection{CC Switch 供应商配置与切换}

CC Switch 3.14.1 已通过 Homebrew Cask 安装，供应商条目保存在 \path{~/.cc-switch/cc-switch.db} 中，按 Codex、Claude Code、Gemini 等工具分类。Codex 分类下当前启用的条目名称为 \verb|pro20x|，该条目携带的配置片段如下（API Key 与 Endpoint 已脱敏）。

\begin{lstlisting}[caption={CC Switch 中启用条目携带的 Codex 配置（脱敏）}]
{
  "auth": { "OPENAI_API_KEY": "<REDACTED>" },
  "config": "model_provider = \"custom\"
             model = \"gpt-5.5\"
             model_reasoning_effort = \"medium\"

             [model_providers.custom]
             name = \"OpenAI\"
             base_url = \"<REDACTED>\"
             wire_api = \"responses\"
             requires_openai_auth = false"
}
\end{lstlisting}

配置与切换的操作步骤为：在 CC Switch 的 Codex 页面点击添加供应商；选择内置预设，或选择自定义后按供应商文档填写名称、API Key、Endpoint URL 和模型，Codex 端点通常以 \verb|/v1| 结尾；保存后点击供应商卡片上的启用；关闭并重新打开 Codex 或终端，使新的供应商和模型列表生效；在会话中用 \verb|/status| 核对当前模型与权限。

切换的生效机制为：CC Switch 将所选条目的 \verb|config| 字段写入 \path{~/.codex/config.toml}，将 \verb|auth| 字段写入 \path{~/.codex/auth.json}。命令行工具在启动时读取配置文件，因此切换后需要重启 Codex CLI 或 Codex App。前述两次 \verb|codex exec| 调用即在当前启用的供应商下完成，返回结果正常，可作为供应商配置生效的验证。

\section{实验总结}

本次实验完成了三部分内容。LaTeX 部分建立了图片规范：结果图由绘图程序直接导出为 PDF 矢量文件，通过 \verb|\includegraphics| 按 \verb|\textwidth| 比例设置尺寸，并配置图题、标签与正文引用。Git 部分完成了仓库创建、克隆、两次本地提交、上传和一次远程更新的完整流程，提交前均通过 \verb|status| 与 \verb|diff| 核对变化范围。Codex 部分先以只读沙盒确认项目理解正确，再在 workspace-write 沙盒下完成编程任务，退出后用 Git 独立核对修改并运行测试。CC Switch 部分核对了供应商条目的存储位置和切换生效机制。

实验过程中形成的几点认识如下。第一，先执行只读任务再执行写入任务，可以在授予写权限前确认工具对项目的理解是否正确。第二，Codex 生成的修改需要通过 \verb|git status| 和 \verb|git diff| 独立核对，工具自身报告的修改范围不作为唯一依据。第三，矢量图应由绘图程序在数据阶段直接导出，重新截图会使前一步的清晰度损失无法恢复。第四，API Key 出现在 CC Switch 界面、\verb|auth.json| 和终端输出中，报告与仓库中的对应位置需要逐一脱敏。

% \bibliographystyle{ieeetr}
% \bibliography{reference}

\end{document}
